%%%%%%%%%%%%%%%%%%%%%%%%%%%%%%%%%%%%%%%%%%%%%%%%%%%%%%%%%%%%%%%%%%%%%%
% Lớp 11 | Đạo hàm | Đạo hàm
% Dạng 1D7?1-2: Số gia hàm số, số gia biến số
% Số câu: 10  (sinh tự động bởi tach_cau.py; từ nay thầy sửa trực tiếp file này)
%%%%%%%%%%%%%%%%%%%%%%%%%%%%%%%%%%%%%%%%%%%%%%%%%%%%%%%%%%%%%%%%%%%%%%

%%%%%%%%%%%%%%%% TRẮC NGHIỆM 4 PHƯƠNG ÁN %%%%%%%%%%%%%%%%

% nguồn: Ok_TracNghiemLop11_Lan1
\begin{ex}%[1D7H1-2]%[ma:042d3161]
	Tỉ số $\dfrac{\Delta y}{\Delta x}$  của hàm số $f ( x)=2 x \cdot ( x -1)$ theo $x$ và $\Delta x$ là
	\choice 
		{$4  x+2 \Delta x+2$}
		{$4  x+2(\Delta x)^2+2$}
		{\True $4  x+2 \Delta x-2$}
		 {$4 x \cdot \Delta x+2(\Delta x)^2-2 \Delta x$}
	\loigiai{ 
		\begin{align*}
			f (x)&=2 x .(  x -1)=2 x ^2-2 x \\
			\Delta y&=f(x+\Delta x)-f(x)  \\
			& =2(x+\Delta x)^2-2 \cdot(x+\Delta x)-\left(2 x^2-2 x\right) \\ 
			& =2 x^2+4 x \cdot \Delta x+2(\Delta x)^2-2 x-2 \Delta  x-2 x^2+2 x \\
			& =4 x \cdot \Delta x+2(\Delta x)^2-2  \Delta x.
		\end{align*}
		Lập tỉ số $\dfrac{\Delta y}{\Delta  x}=4 x+2 \cdot \Delta x-2$.
	 }
\end{ex}

% nguồn: Ok_TracNghiemLop11_Lan1
\begin{ex}%[1D7H1-2]%[ma:f91bccbc]
	Số gia của hàm số $f(x)=x^3$ ứng với $x_0=2$ và $\Delta  x=1$ bằng bao nhiêu?
	\choice 
		{$  -19 $}
		 {$ 7 $}
		{\True $ 19  $}
		{$  -7 $}
	 \loigiai{
		Gọi $\Delta x$ là số gia của đối số và  $\Delta y$ là số gia tương ứng của hàm số. Ta có
		\begin{align*}
			 \Delta y & =f\left(x_0+\Delta x\right)-f\left(x_0\right)\\
			& =\left(x_0+\Delta  x\right)^3-2^3 \\
			& = x_0^3+(\Delta x)^3+3  x_0 \Delta x\left(x_0+\Delta x\right)-8.
		\end{align*}
		Với $x_0=2$ và $\Delta x=1$ thì $\Delta  y=19$.
	} 
\end{ex}

% nguồn: Ok_TracNghiemLop11_Lan1
\begin{ex}%[1D7H1-2]%[ma:53a47f7d]
	Tính  số gia của hàm số $y=\sqrt{2 x+1}$ tại $x _0=1$.
	 \choice
		{$ \dfrac{2 \Delta x}{\sqrt{3+2  \Delta x}-\sqrt{3}} $}
		 {\True $ \dfrac{2 \Delta x}{\sqrt{3+2 \Delta x}+\sqrt{3}} $}
		{$ \dfrac{2 \Delta  x}{\sqrt{3-2 \Delta x}+\sqrt{3}} $}
		 {Đáp án khác}
	 \loigiai{
		Cho $x_0=1$ một số gia $\Delta x$.  Khi đó hàm số nhận một số gia tương ứng là
		\begin{align*}
			\Delta y&=f\left(x_0+\Delta x\right)-f\left(x_0\right)  \\
			& =f(1+\Delta x)-f(1) \\ 
			 & =\sqrt{2(1+\Delta x)+1}-\sqrt{3} \\
			 & =\sqrt{2 \Delta x+3}-\sqrt{3} \\
			& =\dfrac{2 \Delta x}{\sqrt{3+2  \Delta x}+\sqrt{3}}.
		\end{align*}
	 }
\end{ex}

% nguồn: Ok_TracNghiemLop11_Lan1
\begin{ex}%[1D7H1-2]%[ma:1367b761]
	Cho hàm số $f ( x)= x _2- x$, đạo  hàm của hàm số ứng với số gia của đối số $x$ tại $x _0$ là
	 \choice
		{$\lim\limits_{\Delta x \rightarrow 0}\left((\Delta  x)^2+2 x \Delta x-\Delta x\right)$}
		{\True $\lim\limits_{\Delta x \rightarrow 0}(\Delta  x+2 x-1)$}
		 {$\lim\limits_{\Delta x \rightarrow 0}(\Delta x+2 x+1)$}
		{$\lim\limits_{\Delta x \rightarrow 0}\left((\Delta x)^2+2 x  \Delta x+\Delta x\right)$}
	 \loigiai{
		Ta  có
		\begin{align*}
			\Delta y&=\left(x_0+\Delta x\right)^2-\left(x_0+\Delta  x\right)-\left(x_0^2-x_0\right) \\
			& =x_0^2+2 x_0 \Delta x+(\Delta x)^2-x_0-\Delta x-x_0^2+x_0 \\ 
			& =(\Delta x)^2+2  x_0 \Delta x-\Delta x.
		\end{align*}
		Nên 
		\begin{align*}
			f'\left(x_0\right)&=\lim\limits_{\Delta x \rightarrow 0}  \dfrac{\Delta y}{\Delta x} \\
			&  =\lim\limits_{\Delta x \rightarrow 0} \dfrac{(\Delta x)^2+2 x_0 \Delta x-\Delta x}{\Delta x} \\
			&=\lim\limits_{\Delta x \rightarrow 0}\left(\Delta x+2  x_0-1\right).
		\end{align*}
		Vậy $f'(x)=\lim\limits_{\Delta x  \rightarrow 0}(\Delta x+2 x-1)$.
	} 
\end{ex}

% nguồn: Ok_TracNghiemLop11_Lan1
\begin{ex}%[1D7H1-2]%[ma:c3878bc4]
	Cho hàm số $y=f(x)$ xác định và có đạo hàm trên khoảng $(a; b)$ và $x_0 \in(a; b)$.  Đạo hàm của hàm số $f(x)$ tại $x_0$ là
	 \choice
		{\True $f'\left(x_0\right)=\lim\limits_{x\to  x_0}\dfrac{f(x)-f\left(x_0\right)}{x-x_0}$}
		{$f'\left(x_0\right)=\lim\limits_{x\to  x_0}\dfrac{f(x)+f\left(x_0\right)}{x-x_0}$}
		 {$f'\left(x_0\right)=\lim\limits_{x\to x_0}\dfrac{f(x)-f\left(x_0\right)}{x+x_0}$}
		{$f'\left(x_0\right)=\lim\limits_{x\to  x_0}\dfrac{f(x)+f\left(x_0\right)}{x+x_0}$}
	\loigiai{Theo định nghĩa đạo  hàm của hàm số tại một điểm là $f'\left(x_0\right)=\lim\limits_{x\to x_0}\dfrac{f(x)-f\left(x_0\right)}{x-x_0}$.
	} 
\end{ex}

% nguồn: Ok_TracNghiemLop11_Lan1
\begin{ex}%[1D7H1-2]%[ma:66a1e6f7]
	Số gia của hàm số $f(x)=\dfrac{x^2}{2}$ ứng  với số gia $\Delta x$ của đối số $x$ tại $x_0=-1$ là
	\choice 
		{\True $\dfrac{1}{2}(\Delta x)^2-\Delta  x$}
		{$\dfrac{1}{2}\left[(\Delta x)^2-\Delta  x\right]$}
		 {$\dfrac{1}{2}\left[(\Delta x)^2+\Delta x\right]$}
		{$\dfrac{1}{2}(\Delta  x)^2+\Delta x$}
	\loigiai{ 
		Với số gia $\Delta x$ của  đối số $x$ tại $x_0=-1$, ta có
		\begin{align*}
			 \Delta y&=\dfrac{(-1+\Delta x)^2}{2}-\dfrac{1}{2} \\
			&=\dfrac{1+(\Delta x)^2-2  \Delta x}{2}-\dfrac{1}{2} \\
			&=\dfrac{1}{2}(\Delta x)^2-\Delta x. 
		\end{align*}
	} 
\end{ex}

% nguồn: Ok_TracNghiemLop11_Lan1
\begin{ex}%[1D7H1-2]%[ma:3d5924dd]
	Cho hàm số $y=f(x)$ có đạo hàm tại $x_0$ là $f'(x_0)$. Khẳng định nào sau  đây là \textbf{sai}?
	\choice 
		{\True  $f'(x_0)=\lim\limits_{x\to x_0}\dfrac{f(x+x_0)-f(x_0)}{x-x_0}$}
		{$f'(x_0)=\lim\limits_{x\to \Delta x}\dfrac{f(x_0+\Delta  x)-f(x_0)}{\Delta x}$}
		{$f'(x_0)=\lim\limits_{x\to  x_0}\dfrac{f(x)-f(x_0)}{x-x_0}$}
		{$f'(x_0)=\lim\limits_{h\to 0}\dfrac{f(h+x_0)-f(x_0)}{h}$} 
	\loigiai{ 
		Theo định nghĩa đạo  hàm của hàm số tại một điểm.\\
		Mệnh đề  $f'(x_0)=\lim\limits_{x\to x_0}\dfrac{f(x+x_0)-f(x_0)}{x-x_0}$ là mệnh đề sai
	 }
\end{ex}

% nguồn: Ok_TracNghiemLop11_Lan1
\begin{ex}%[1D7Y1-2]%[ma:74fc45ea]
	Số gia $\Delta y$ của hàm số  $y = x^2 + 2x - 5$ tại điểm $x_0 = 1$ là
	 \choice
		{$(\Delta x)^2 -  4\Delta x$}
		{\True $\left( \Delta x \right)^2  + 4\Delta x$}
		{$\left(\Delta x  \right)^2 - 2\Delta x$}
		{$\left( \Delta x  \right)^2 + 2\Delta x - 5$}
	 \loigiai{
		Ta có $$\begin{aligned} \Delta y&=f(1+\Delta x)-f(1)=\left(1+\Delta x \right)^2+2(1+\Delta x)-5+2 \\&= \left(\Delta x \right)^2+4\Delta x.
	\end{aligned} $$
 }
\end{ex}


%%%%%%%%%%%%%%%% ĐÚNG/SAI %%%%%%%%%%%%%%%%

% nguồn: Ok_dungsai_lop11_Lan1
\begin{ex}%[1D7N1-2]%[ma:e59f24cc]
 Cho hàm số  $f(x)=x^2+2x$, có $\Delta x$ là số gia của đối số tại $x=1$, $\Delta y$ là số gia tương ứng của hàm số. Khi đó, các khẳng định sau đúng hay sai?
  \choiceTF
   {$\dfrac{\Delta y}{\Delta x} = \Delta x+2$}
   {\True $\Delta y = (\Delta x)^2+4 \Delta x$}
  {$\Delta y - 4\Delta x = (\Delta x)^2 -  2 \Delta x-3$}
  {\True $2\Delta y - 6\Delta x  = 2(\Delta x)^2 + 2 \Delta x$}
 \loigiai{ 
  Ta có 
  $$\Delta y=f(1+\Delta  x)-f(1)=(1+\Delta x)^2+2(1+\Delta x)-(1+2)=(\Delta x)^2+4 \Delta x.$$
  $\Rightarrow 2\Delta y - 6\Delta x = 2(\Delta x)^2 +  2 \Delta x$.
  }
\end{ex}

% nguồn: Ok_dungsai_lop11_Lan1
\begin{ex}%[1D7H1-2]%[ma:accb3735]
 Các khẳng định dưới  đây đúng hay sai?
  \choiceTF
  {\True Số gia của hàm số $y=x^3+x^2+1$ tại điểm $x_0$ ứng với số  gia $\Delta x=1$ là $\Delta y=3x_0^2+5x_0+3$}
  {Số gia của hàm  số $y=\dfrac{x^2}{2}$ tại điểm $x_0=-1$ ứng với số gia $\Delta x$ là $\Delta y=\dfrac{1}{2}[(\Delta x)^2-\Delta x]$}
  {Gọi $\Delta x$ là số gia của đối số tại $x$ và $\Delta y$ là số gia tương ứng của hàm số $y=x^3+1$. Khi đó $\dfrac{\Delta y}{\Delta x} = 3x^2-3x\cdot \Delta x+(\Delta x)^3$} 
  {\True Tỉ số $\dfrac{\Delta y}{\Delta x}$  của hàm số $f(x)=2x(x-1)$ theo $x$ và $\Delta x$ là $4x+2 \cdot \Delta x-2$}
  \loigiai{
  \begin{itemchoice}
   \itemch $\Delta y=f(x_0+\Delta x)-f(x_0)=f(x_0+1)-f(x_0)=[(x_0+1)^3+(x_0+1)^2+1]-(x_0^3+x_0^2+1) =3x_0^2+5x_0+2$. 
    \itemch $\Delta y=f(-1+\Delta x)-f(-1)=\dfrac{(-1+\Delta x)^2}{2}-\dfrac{1}{2}=\dfrac{1}{2}[(-1+\Delta x)^2-1]=\dfrac{1}{2}(\Delta x)^2-\Delta x$.
   \itemch Ta có: $\begin{aligned}
   \Delta y&=f(x+\Delta  x)-f(x)\\
   &=(x+\Delta x)^3+1-(x^3+1)\\ 
   &=3x^2\cdot  \Delta x+3x\cdot \Delta ^2x+\Delta ^3x\\
    &=\Delta x(3x^2+3x\cdot \Delta x+\Delta ^2x).
  \end{aligned}$\\
   $\Rightarrow \dfrac{\Delta y}{\Delta x}=3x^2+3x\cdot \Delta x+\Delta ^2x=3x^2+3x\cdot \Delta x+(\Delta x)^2$.
  \itemch Ta có:  $f(x)=2x(x-1)=2x^2-2x$.\\
  $\begin{aligned}
  \Delta y&=f(x+\Delta  x)-f(x)\\
  & =2(x+\Delta x)^2-2 \cdot(x+\Delta x)-\left(2 x^2-2  x\right) \\
  & =2 x^2+4  x \cdot \Delta x+2(\Delta x)^2-2 x-2 \Delta x-2 x^2+2 x \\
  & =4 x \cdot \Delta x+2(\Delta x)^2-2 \Delta  x.
 \end{aligned}$\\
  $\Rightarrow \dfrac{\Delta y}{\Delta x}=4x+2 \cdot \Delta x-2$.
 \end{itemchoice}
 }
\end{ex}
